\tinysidebar{\debug{exercises/{test/question.tex}}}
\begin{ex}
A brute force way to compute GCD is this:

\begin{quote}
Given positive integer m and n, run d from 1 to the smaller of m and n,
check if d is a divisor of both m and n. Compute the largest of such d.
\end{quote}

Here's the pseudocode:

\begin{consolethree}
int max_d;
for d = 1, 2, 3, ..., min(m,n):
    if d divides both m and n:
        if d < max_d:
            max_d = d
\end{consolethree}

The max\_d is the GCD of m,n. (Why does this running max computation not
require an initialization of max\_d?) Since d keeps increasing, the
above is really the same as

\begin{consolethree}
int max_d;
for d = 1, 2, 3, ..., min(m,n):
    if d divides both m and n:
        max_d = d
\end{consolethree}

I'll let you think about this \ldots{} this is a better pseudocode:

\begin{consolethree}
int max_d;
for d = min(m,n), min(m,n) - 1, min(m,n) - 2, ..., 1:
    if d divides both m and n:
        max_d = d and break the loop
\end{consolethree}

Now you have two ways to compute GCD: by the earlier recursion or using
this loop. Try to compute the GCD of several pairs of m,n and decide
which algorithm is faster. (Hint: Euclid is smart.)
\end{ex}
